% engines: lualatex, pdflatex
% expect-warnings: 0
% expect-errors: 0
\documentclass{article}
\usepackage[margin=1cm]{geometry}
\usepackage{../../latex/pgf-go}
\graphicspath{{images/}}
\newcommand\case[1]{\par\medskip\noindent\textbf{#1}\par\nopagebreak}
\pagestyle{empty}
\pgfgodefinemark{star}[image = star.pdf]
\pgfgodefinemark{big star}[image = star.pdf, width = 14pt]
\pgfgodefinemark{diamond}{%
	\pgfpathmoveto{\pgfpointxy{0}{0.7}}\pgfpathlineto{\pgfpointxy{0.7}{0}}%
	\pgfpathlineto{\pgfpointxy{0}{-0.7}}\pgfpathlineto{\pgfpointxy{-0.7}{0}}%
	\pgfpathclose\pgfusepath{fill, stroke}}
\pgfgodefinemark{bar}{\draw (0,0) circle (0.5); \draw[pgfgomarkfill] (-0.5,0) -- (0.5,0);}
\begin{document}

\case{Stone images (PNG with transparency) with a shadow of the same shape; two white images chosen by position; an oval stone; a prisoner keeps its opacity; marks on image stones}
\begin{goban}[9, label = {at = {south, west}}]
	\stones[b, image = {file = stone-black.png}, drop shadow = {true, angle = -45, offset = 2pt}]{A1, B2, C3, D4}
	\stones[w, image = {file = {stone-white.png, stone-white-2.png}}, drop shadow = {true, angle = -45, offset = 2pt}]{E5, F5, G5, H5, E6, F6, G6, H6}
	\stones[b, image = {file = stone-oval.png, width = 20pt}, drop shadow = {true, offset = 2pt, angle = 90}]{C7}
	\prisoners[w, image = {file = stone-white.png}]{A9}
	\marks[triangle]{A1, E5}
\end{goban}
\begin{goban}[5, scale = 2, partial = {from = A1 to C3}]
	\stones[b, image = {file = stone-black.png}, drop shadow = {true, angle = -45, offset = 2pt}]{A1, B2}
	\stones[w, image = {file = stone-white.png, width = 8pt}]{C3}
\end{goban}

\case{Rotated images: random (one angle per point, the shadow keeps its direction), a fixed angle, and random as a player key}
\begin{goban}[5]
	\stones[b, image = {file = stone-oval.png, width = 18pt, rotate = random}, drop shadow = {true, angle = -45, offset = 2pt}]{A1, B1, C1, D1, E1, A2, B2, C2, D2, E2}
	\stones[w, image = {file = stone-oval.png, width = 18pt, rotate = 45}, drop shadow = {true, angle = -45, offset = 2pt}]{A4, B4}
	\stones[w, image = {file = stone-oval.png, width = 18pt, rotate = 0}]{D4}
\end{goban}
{\pgfgoset{players = {black = {stone = {image = {file = stone-oval.png, width = 18pt, rotate = random}}}}}
\begin{goban}[5]
	\moves{A1, B2, C3, D4}
\end{goban}}

\case{Rotation ranges: 80:100 (almost vertical), -10:10 (almost horizontal), and spaces around the colon}
\begin{goban}[5]
	\stones[b, image = {file = stone-oval.png, width = 18pt, rotate = 80:100}]{A1, B1, C1, D1, E1}
	\stones[w, image = {file = stone-oval.png, width = 18pt, rotate = -10:10}]{A3, B3, C3, D3, E3}
	\stones[b, image = {file = stone-oval.png, width = 18pt, rotate = 80 : 100}]{A5, B5, C5}
\end{goban}

\case{Player keys: moves use the image; image = false goes back to circles}
\pgfgoset{players = {black = {stone = {image = {file = stone-black.png}}}, white = {stone = {image = {file = stone-white.png}}}}}
\begin{goban}[5]\moves{A1, B1, C1, D1}\stones[b, image = false]{A3}\end{goban}
\pgfgoset{players = {black = {stone = {image = false}}, white = {stone = {image = false}}}}
\begin{goban}[5]\moves{A1, B1}\end{goban}

\case{Background images: cover, stretch (rounded corners, partial), tile (partial); marks on empty points show the image under them; dim}
\begin{goban}[9, background = {fill = {image = {file = wood.jpg}}}]
	\stones{C3}\stones[w]{D4}
	\marks[circle]{A1, E5, C3}\marks[dim]{F6, D4}
\end{goban}
\begin{goban}[9, background = {fill = {image = {file = wood.jpg, fit = stretch}}, rounded corners = 6pt}, partial = {from = A1 to E5}]
	\marks[circle]{A1, E5}
\end{goban}
\begin{goban}[9, background = {fill = {image = {file = tile.png, fit = tile, width = 1cm}}}, partial = {from = D4 to H8}]
	\marks[triangle]{E5}
\end{goban}
\begin{goban}[5, background = {fill = {image = {file = wood.jpg, false}}}]\marks[circle]{C3}\end{goban}

\case{Tiles from corners and sides of the whole board and from the centre}
\begin{goban}[5, background = {fill = {image = {file = tile.png, fit = tile, width = 0.9cm, from = south east}}}]\end{goban}
\begin{goban}[5, background = {fill = {image = {file = tile.png, fit = tile, width = 0.9cm, from = north west}}}]\end{goban}
\begin{goban}[5, background = {fill = {image = {file = tile.png, fit = tile, width = 0.9cm, from = center}}}]\end{goban}
\begin{goban}[5, background = {fill = {image = {file = tile.png, fit = tile, width = 0.9cm, from = north}}}]\end{goban}
\begin{goban}[5, background = {fill = {image = {file = tile.png, fit = tile, width = 0.9cm, from = west}}}]\end{goban}
\begin{goban}[9, background = {fill = {image = {file = tile.png, fit = tile, width = 0.9cm, from = north east}}}, partial = {from = D4 to H8}]\end{goban}

\case{Grid and hoshi images; the grid image is clipped by partial}
\begin{goban}[9, grid = {image = {file = grid.png}}, hoshi = {image = {file = hoshi.png}, radius = 3pt}, background = {fill = {image = {file = tile.png, fit = tile}}}]
	\stones{C3}
\end{goban}
\begin{goban}[9, grid = {image = {file = grid.png}}, partial = {from = C3 to G7}]\end{goban}

\case{Grid image as tiles: one per square of the grid (also with partial and with stones), a 4x4 and a 2x4 mesh; a 3x3 mesh on the background}
\begin{goban}[9, grid = {image = {file = cell.png, fit = tile}}, hoshi = false, background = {fill = {color = black!70}}]
	\stones{C3}\stones[w]{D4}
\end{goban}
\begin{goban}[9, grid = {image = {file = cell.png, fit = tile}}, partial = {from = C3 to G7}]\end{goban}
\begin{goban}[5, grid = {image = {file = cell.png, fit = tile, false}}]\end{goban}
\begin{goban}[9, grid = {image = {file = cell.png, fit = tile, tiles = 4}}, hoshi = false]\end{goban}
\begin{goban}[9, grid = {image = {file = cell.png, fit = tile, tiles = 2x4}}, partial = {from = C3 to G7}]\end{goban}
\begin{goban}[9, background = {fill = {image = {file = cell.png, fit = tile, tiles = 3x3}}}]\end{goban}

\case{Custom marks: PDF image (default width and 14pt), pgf code with the mark colours, TikZ code; on black and white stones; as a territory mark}
\begin{goban}[7]
	\stones{B2}\stones[w]{D2}
	\marks[star]{A1, B2}\marks[big star]{C1}
	\marks[diamond]{A3}\marks[b, diamond]{B2}\marks[w, diamond]{D2}
	\marks[bar]{A5}\marks[b, bar]{E5}
	\marks[w, fill = {true, color = red}, bar]{G5}
	\territory[b, mark = star]{F1, G1}
\end{goban}
\end{document}
